% ==============================================================================
% apendices/apendice_a_checklist.tex
% Apêndice A — Checklist de Comissionamento e Implantação em Bancada
% ==============================================================================

\chapter{Checklist de Comissionamento e Implantação em Bancada}
\label{ap:checklist}

Este apêndice consolida a sequência operacional recomendada para montagem, configuração inicial e verificação pré-ensaio da bancada envolvendo o controlador WEG PLC500, o ambiente CODESYS V3.5, o banco PostgreSQL e o fluxo no Node-RED.

\section{Etapa 1 -- Montagem Elétrica e Verificação de Hardware}
\begin{enumerate}
  \item Confirmar alimentação regulada de 24\,Vcc no WEG PLC500 e proteção por disjuntor/fusível adequado à bancada didática.
  \item Conectar as botoeiras NA em \texttt{DI0} (\texttt{\%IX0.0}), \texttt{DI1} (\texttt{\%IX0.1}) e \texttt{DI2} (\texttt{\%IX0.2}) em configuração PNP (comutando $+24\,\text{Vcc}$ para a entrada).
  \item Conectar a botoeira de emergência NF em \texttt{DI3} (\texttt{\%IX0.3}), confirmando presença de nível lógico alto (LED indicador da entrada aceso no CLP) em condição de repouso.
  \item Conectar os quatro sinalizadores LED de 24\,Vcc (ou LEDs discretos com resistores limitadores de $1{,}8\text{--}2{,}2\,\text{k}\Omega$) nas saídas \texttt{DO0}--\texttt{DO3} (\texttt{\%QX0.0}--\texttt{\%QX0.3}) referenciados ao $0\,\text{Vcc}$.
  \item Conectar o cabo Ethernet de engenharia à porta \texttt{ETH1} e o cabo da sub-rede Modbus dedicada à porta \texttt{ETH2}.
\end{enumerate}

\section{Etapa 2 -- Configuração do WEG PLC500 no CODESYS V3.5}
\begin{enumerate}
  \item Configurar o IP estático \texttt{192.168.20.10/24} na interface \texttt{ETH1} e \texttt{192.168.10.10/24} na interface \texttt{ETH2}.
  \item Inserir \texttt{Ethernet\_ETH2} e o dispositivo filho \texttt{ModbusTCP\_Server\_Device} (porta \texttt{502}, \textit{Unit ID} \texttt{1}) na árvore de dispositivos.
  \item Mapear integralmente os endereços base zero conforme a Tabela~\ref{tab:mapa_modbus}, habilitando a opção de atualização contínua das variáveis no ciclo de \texttt{MainTask} ($10\,\text{ms}$).
  \item Compilar sem erros, realizar \textit{download} de \texttt{PLC\_PRG} (Quadro~\ref{quadro:codigo_st}), colocar o CLP em modo \texttt{RUN} e verificar \texttt{uiVersao\_Mapa = 100} no registrador \texttt{HR[3]}.
\end{enumerate}

\section{Etapa 3 -- Implantação do PostgreSQL e do Fluxo Node-RED}
\begin{enumerate}
  \item No servidor \texttt{192.168.10.20}, verificar conectividade ICMP/TCP na porta \texttt{502} com \texttt{192.168.10.10} e confirmar regra de firewall restringindo a porta 502 apenas ao host do Node-RED.
  \item Criar a base de dados PostgreSQL dedicada e executar o script DDL do Quadro~\ref{quadro:ddl_sql}, validando a criação das quatro tabelas e do índice \texttt{idx\_tel\_capturado}.
  \item No gerenciador de paleta do Node-RED (\textit{Manage Palette}), instalar os pacotes \texttt{node-red-contrib-modbus} e \texttt{node-red-contrib-postgresql}, configurando o nó \texttt{Modbus-Client} (\texttt{192.168.10.10:502}, \textit{Unit ID}~\texttt{1}, \textit{timeout}~$1500\,\text{ms}$).
  \item Estruturar o fluxo conforme a Figura~\ref{fig:fluxo-pipeline-nodered}, inserindo o código JavaScript do Quadro~\ref{quadro:codigo_js} no nó \texttt{function} de análise temporal e o contrato do Quadro~\ref{quadro:json_schema} no nó validador JSON após o nó \texttt{http request}.
  \item Ativar o modo remoto (\texttt{Coil[4] = TRUE}), confirmar o acendimento do indicador \texttt{L4\_StatusOper} (\texttt{DO3}) e iniciar sequencialmente a execução dos casos de teste \texttt{CT-01} a \texttt{CT-10} do Quadro~\ref{quadro:casos_teste}.
\end{enumerate}
